\begin{abstract}
Prediction Shift Backdoor Detection (PSBD) flags a poisoned input by how little its prediction moves when dropout is switched on at inference: a backdoored input rides a shortcut that survives the perturbation, a clean input does not. The method was designed and validated on ResNets. We ask whether it transfers to Vision Transformers, and find that the answer depends almost entirely on where the perturbation is injected. Ported literally, with dropout after the residual add as in the original paper, PSBD reaches a mean AUROC of 0.832 over 65 backdoored ViT-B/16 checkpoints spanning 4 datasets, 7 attacks and poison rates from 1\% to 10\%, with 12 of the 65 cells inverted. Masking whole tokens at the input of each attention block instead reaches 0.935, a paired gain of +0.103 (95\% bootstrap interval +0.048 to +0.160), largest where detection is hardest: +0.110 on the hard attacks and +0.161 at 1\% poisoning. The choice of operator at the same site matters as much as the site: Gaussian noise at the attention input loses 0.206 to token masking, while the same noise at the MLP input wins by 0.056, and we trace both to how LayerNorm absorbs additive noise but not masking. Following the backdoor through the network explains why the placement works. Across attacks the backdoor is 1 direction in the residual stream, written in the last third of the network and routed from the trigger's patch tokens to the class token through attention, so removing tokens before that routing removes the shortcut. The prediction shift phenomenon the original paper attributes to neuron bias exists on ViT only for local triggers on few-class datasets and is absent for global triggers, while detection works in both regimes, so the mechanism on transformers is decision-margin estimation with routing as the special case. The same placement transfers to Swin-S with a gain of +0.091 over the published one. We also report what does not hold: a single probe is defeated by an adaptive attacker that trains against it, and only a union over probes restores detection, and PSBD does poorly on all-to-all attacks. All 65 cells, the 18-placement basis and every number are generated from a coverage ledger and released with the code.
\end{abstract}
